\subsection{একটি কলেজে ভিন্ন ধরনের ডেটা আদান-প্রদান}
\label{ch2:c02-14}
কলেজের সেন্সর অনিয়মিত ছোট বার্তা পাঠায়, অফিস বড় ফাইল আদান-প্রদান করে এবং শিক্ষক সরাসরি অডিও ক্লাস নেন। কাজগুলোর সময়, ডেটার পরিমাণ ও প্রাপকের চাহিদা আলাদা। তাই একই প্রেরণব্যবস্থা সব কাজে সমান সুবিধা দেয় না।
\chTwoPair{\textbf{সেন্সর ও বড় File}\par সেন্সরের ছোট ক্যারেক্টারভিত্তিক বার্তায় অ্যাসিনক্রোনাস সংযোগ বিবেচনা করা যায়। বড় ফাইলের ডেটা বাফারে জমিয়ে সিনক্রোনাস ব্লকে পাঠালে বাড়তি তথ্য প্রতিটি ক্যারেক্টারের বদলে পুরো ব্লকের জন্য ব্যবহৃত হয়। Efficiency ও পুনঃপ্রেরণের হিসাব অংশ \ref{ch2:c02-07}-এ আছে।}{\textbf{Live Class}\par শব্দের নমুনা নিয়মিত সময়ে পৌঁছানো দরকার। ব্যবস্থার সময়সূচি জেনে আইসোক্রোনাস প্রেরণের প্রয়োজন বিচার করতে হয়। শুধু অ্যাপের নাম থেকে প্রোটোকল নিশ্চিত করা যায় না। একই সময়ে কথা বলা ও শোনা সম্ভব হলে যোগাযোগ ফুল-ডুপ্লেক্স। শুধু একদিকে বক্তৃতা পাঠানো হলে তা সিমপ্লেক্সও হতে পারে।}
মাল্টিকাস্ট সমর্থিত ব্যবস্থায় নিবন্ধিত দলীয় ঠিকানায় ডেটা পাঠালে নির্বাচিত শিক্ষার্থীরা একই প্রবাহ পেতে পারে। একজনের ফাইল ডাউনলোড এবং দলের সরাসরি ক্লাসে প্রাপক আলাদা। তবে শুধু প্রাপক সংখ্যা থেকে যোগাযোগের দিক বা সময়ের নিয়ম জানা যায় না। একই কাজে সিরিয়াল পথ, ব্লকভিত্তিক সময় সমন্বয়, ফুল-ডুপ্লেক্স দিক ও ইউনিকাস্ট প্রাপক থাকতে পারে। শ্রেণিগুলো যোগাযোগের পৃথক দিক বোঝায় (অংশ \ref{ch2:c02-01})।
